\section{Dataset Construction}
\label{sec:app_data}
% =============================================================================

This section documents how the Stage~1 (projector alignment) and Stage~2 (task-specific RL) training corpora are assembled. The corresponding test-time disjointness guarantees---FEN-level non-overlap between every LLAMIA-Bench test split and the training pools described below, including the heuristic used for Agadmator-2K---are stated once in \cref{sec:app_bench_contamination} and are not repeated here.

% =============================================================================
\subsection{Stage~1: Projector Alignment Data}
\label{sec:app_data_stage1}

% Note: train/test disjointness for projector data is covered in \cref{sec:app_bench_contamination}; only construction is described here.

Stage~1 trains the LatentBridge projector $H_\varphi$ on state--policy pairs from the Lc0-BT4 forward pass. We construct the dataset as follows:

\paragraph{Source.} We sample 5M positions from the Lichess evaluation database (January 2013 -- December 2024), filtering for standard-time-control games between rated players ($\geq$1200 Elo). Positions are sampled uniformly across game phases (opening: moves 1--15, middlegame: moves 16--35, endgame: moves 36+) to prevent phase bias.

\paragraph{Label generation.} For each position, we run a single BT4 forward pass (no MCTS search) to obtain the raw policy distribution $\pi_\text{BT4}(s)$, the value head output $V(s)$, and the penultimate-layer activations $\vh_s \in \R^{1024}$. The training target is the top-1 move from the policy head, formatted as either UCI or SAN notation (70\%\,/\,30\%). The activation $\vh_s$ is the input to the projector.

\paragraph{Prompt diversity.} Each position is paired with one of four question types (Section~\ref{sec:app_prompts_stage1}): position evaluation, principal variation, legal moves, or brief description. Question types are sampled uniformly. This diversity prevents the projector from overfitting to a single output format.

\paragraph{Split.} The 5M positions are split by game ID (not by position) to prevent train/test leakage: 4.5M training, 250K validation, 250K held-out test. No game appears in more than one split.

% =============================================================================
\subsection{Stage~2: Task-Specific RL Data}
\label{sec:app_data_stage2}

Stage~2 uses DAPO rollouts on task-specific prompts. The training data for RL totals ${\sim}$850K examples across all tasks:

\paragraph{Behavior Cloning.} 500K positions from Lichess games, stratified by player Elo (100-point bins from 1100 to 2600). Each position is paired with the move actually played by the human player. The reward signal is based on rank within the engine's top-3 moves: the model receives reward 1.0 for a top-1 match, 0.5 for top-2, 0.25 for top-3, and 0 otherwise. Top-1 exact match as a reward collapsed training; rank within the top-3 provides a denser, monotone signal.

\paragraph{Puzzle Understanding.} 200K puzzles from the Lichess puzzle database, each annotated with difficulty rating and popularity score. The reward is a scaled negative absolute error between LLAMIA's prediction and the ground truth.

\paragraph{Move Annotation.} 100K annotated positions drawn from 90K games in the GameKnot~\cite{Gameknot} and Lichess annotation corpora (multiple annotations per game). The reward is a G-eval score (GPT-4o judge) comparing LLAMIA's annotation to the reference.

\paragraph{Game Commentary.} 50K annotated game segments (15--30 moves each) from grandmaster commentary databases, chess books transcribed to PGN, and Lichess studies with annotations. The reward combines a G-eval score for commentary quality with BLEU-2 against reference commentaries.

% =============================================================================
